\documentclass[12pt]{article}
\usepackage[T1]{fontenc}
\usepackage[utf8]{inputenc}
\usepackage{lmodern}
\usepackage{textcomp}
\usepackage{enumitem}
\usepackage{geometry}
\geometry{margin=1in}
\begin{document}
\begin{center}
Answer Key - History - Undergraduate Level Assessment 20 \
\small (Answers are ordered exactly as in the question paper.)
\end{center}

\section*{Multiple Choice}
\begin{enumerate}[label=\arabic*.]
\item \textbf{Answer:} A
\\ \textit{Options:} A) in the Mesolithic, goods traveled less far and fewer goods were exchanged; B) in the Mesolithic, goods traveled farther but the same number of goods were exchanged; C) in the Mesolithic, goods traveled less far but the same number of goods were exchanged; D) in the Mesolithic, goods traveled the same distance but more goods were exchanged; E) in the Mesolithic, goods traveled farther but fewer goods were exchanged
\\ \textit{Note:} The correct answer highlights that during the Mesolithic period, the development of more localized economies and the reliance on regional resources led to shorter trade distances and a decrease in the volume of goods exchanged compared to the more extensive trade networks of the Upper Paleolithic, which were characterized by long-distance exchanges of valuable materials.
\item \textbf{Answer:} A
\\ \textit{Options:} A) establishment of permanent settlements; B) development of agriculture; C) domestication of animals; D) creation of written language; E) symbolic expression through the production of art
\\ \textit{Note:} The establishment of permanent settlements is considered the most ancient characteristic of modern humans because it marks a significant shift in human lifestyle from nomadic hunter-gatherer groups to sedentary agricultural societies, which began well before the Late Stone Age and the emergence of anatomically modern humans.
\item \textbf{Answer:} B
\\ \textit{Options:} A) Inca; AD 200; B) Moche; AD 400; C) Olmec; AD 100; D) Maya; AD 400; E) Zapotec; AD 500
\\ \textit{Note:} The Moche civilization, known for its advanced agricultural practices and impressive architectural achievements, emerged in northern Peru and reached its peak around AD 400, making it one of the earliest and most notable kingdoms in South America.
\item \textbf{Answer:} C
\\ \textit{Options:} A) 3100 B.P.; B) 4100 B.P.; C) 5100 B.P.; D) 6100 B.P.
\\ \textit{Note:} The first pharaohs emerged around 5100 years before present (B.P.) as part of the early dynastic period of ancient Egypt, marking the establishment of centralized rule and the formation of a unified state.
\item \textbf{Answer:} C
\\ \textit{Options:} A) had to pay a high tax in goods and labor to the local governor.; B) lived under constant threat of being sold into slavery.; C) were mostly left alone and did fairly well for themselves.; D) were under the direct control of the high priest, who demanded regular sacrifices.; E) were heavily policed by the king's guards to prevent any form of uprising.
\\ \textit{Note:} The commoners of the rural Aztec villages of Capilco and Cuexcomate were able to sustain their livelihoods and manage their agricultural practices independently due to the relatively decentralized nature of governance in these regions, which allowed them a degree of autonomy and stability.
\end{enumerate}

\section*{True / False}
\begin{enumerate}[label=\arabic*.]
\setcounter{enumi}{5}
\item \textbf{Answer:} True
\\ \textit{Options:} A) True; B) False
\\ \textit{Note:} According to the passage: by supposing there was some supernatural intervention on the path leading to humans
\item \textbf{Answer:} True
\\ \textit{Options:} A) True; B) False
\\ \textit{Note:} According to the passage: France had few counterbalancing European successes.
\item \textbf{Answer:} True
\\ \textit{Options:} A) True; B) False
\\ \textit{Note:} According to the passage: foreign officials
\item \textbf{Answer:} True
\\ \textit{Options:} A) True; B) False
\\ \textit{Note:} According to the passage: people had to hire scribes
\item \textbf{Answer:} True
\\ \textit{Options:} A) True; B) False
\\ \textit{Note:} According to the passage: molten
\end{enumerate}

\section*{Long Form Response}
\begin{enumerate}[label=\arabic*.]
\setcounter{enumi}{10}
\item \textbf{Answer:} Anne Hathaway
\\ \textit{Note:} Expected answer: Anne Hathaway
\item \textbf{Answer:} 500 million
\\ \textit{Note:} Expected answer: 500 million
\end{enumerate}

\vfill
\noindent\textit{End of Answer Key}
\end{document}